% =========================================================
% ESTILOS Y COMANDOS PROPIOS DEL INFORME
% Complementa a pucv_inf_2024.sty sin modificarlo.
% =========================================================

% ---------------------------------------------------------
% Rótulos en español
% Se reaplican después de babel para conservar los títulos de
% índices definidos por pucv_inf_2024.sty.
% ---------------------------------------------------------

\AtBeginDocument{%
    \renewcommand{\figurename}{Figura}%
    \renewcommand{\tablename}{Tabla}%
    \renewcommand{\nomname}{\uppercase{Lista de Símbolos}}%
    \renewcommand{\contentsname}{\uppercase{Índice General}}%
    \renewcommand{\listfigurename}{\uppercase{Lista de Figuras}}%
    \renewcommand{\listtablename}{\uppercase{Lista de Tablas}}%
    \renewcommand{\chaptername}{Capítulo}%
    \renewcommand{\appendixname}{Anexo}%
}

% ---------------------------------------------------------
% Bibliografía en español
% Fuerza los rótulos APA en español («Recuperado el ... de»,
% «2.ª ed.», «y» en lugar de «and») sin depender de babel.
% ---------------------------------------------------------

\ExecuteBibliographyOptions{language = spanish}
\DeclareLanguageMapping{spanish}{spanish-apa}

% ---------------------------------------------------------
% Colores auxiliares
% ---------------------------------------------------------

\definecolor{notacolor}{HTML}{4472c4}   % azul institucional (igual a titlecolor)
\definecolor{notafondo}{HTML}{EEF2FA}
\definecolor{pendientecolor}{HTML}{B85C00}
\definecolor{pendientefondo}{HTML}{FDF3E7}
\definecolor{filacabecera}{HTML}{DCE6F5}

% ---------------------------------------------------------
% Cajas de nota y de tareas pendientes
% (equivalen a las citas en bloque y a los bloques
%  "POR COMPLETAR" del documento original en Markdown)
% ---------------------------------------------------------

\newtcolorbox{nota}[1][]{%
    enhanced,
    breakable,
    colback = notafondo,
    colframe = notacolor,
    boxrule = 0pt,
    leftrule = 3pt,
    arc = 0pt,
    outer arc = 0pt,
    left = 8pt, right = 8pt, top = 6pt, bottom = 6pt,
    fontupper = \small,
    parbox = false,
    #1
}

\newtcolorbox{pendiente}[1][Por completar]{%
    enhanced,
    breakable,
    colback = pendientefondo,
    colframe = pendientecolor,
    boxrule = 0.6pt,
    arc = 2pt,
    left = 8pt, right = 8pt, top = 6pt, bottom = 6pt,
    fontupper = \small,
    parbox = false,
    title = {\footnotesize\bfseries\MakeUppercase{#1}},
    coltitle = white,
    colbacktitle = pendientecolor
}

% Marca breve, para usar dentro de una línea de texto
\newcommand{\porcompletarinline}[1]{%
    \textcolor{pendientecolor}{\textbf{[por completar: #1]}}%
}

% ---------------------------------------------------------
% Estilo común de las tablas largas
% ---------------------------------------------------------

\newcommand{\estilotabla}{%
    \footnotesize
    \setlength{\parskip}{0pt}%
    \setlength{\parindent}{0pt}%
    \setlength{\tabcolsep}{4pt}%
    \renewcommand{\arraystretch}{1.25}%
}

\newcommand{\estilotablachica}{%
    \estilotabla
    \scriptsize
}

% Separación antes y después de las longtable
\setlength{\LTpre}{8pt}
\setlength{\LTpost}{8pt}

% Columnas de párrafo centradas verticalmente arriba
\newcolumntype{L}[1]{>{\raggedright\arraybackslash}p{#1}}
\newcolumntype{C}[1]{>{\centering\arraybackslash}p{#1}}

% ---------------------------------------------------------
% Identificadores de código, clases y estados
% ---------------------------------------------------------

\newcommand{\clase}[1]{\texttt{#1}}
\newcommand{\estado}[1]{\texttt{#1}}
\newcommand{\este}[1]{\textit{#1}}

% Estereotipos UML
\newcommand{\ster}[1]{\texttt{\guillemotleft #1\guillemotright}}

% ---------------------------------------------------------
% Figuras de diagramas: comando abreviado
% #1 ancho relativo · #2 archivo · #3 título breve (Lista de Figuras)
% #4 leyenda completa · #5 etiqueta
% ---------------------------------------------------------

\newcommand{\figuradiagrama}[5]{%
    \begin{figure}[H]
        \centering
        \includegraphics[width=#1\textwidth]{#2}
        \caption[#3]{#4}
        \label{#5}
    \end{figure}%
}

% ---------------------------------------------------------
% Formato de los anexos: "Anexo A." en lugar de "A."
% ---------------------------------------------------------

\newcommand{\formatoanexos}{%
    \renewcommand{\thechapter}{Anexo~\Alph{chapter}}%
    \setcounter{figure}{0}%
    \setcounter{table}{0}%
}

% ---------------------------------------------------------
% Ajustes de listas y espaciado dentro de entornos
% ---------------------------------------------------------

\setlist[itemize]{leftmargin=1.2cm, itemsep=2pt, topsep=4pt}
\setlist[enumerate]{leftmargin=1.2cm, itemsep=2pt, topsep=4pt}

% Evita que hyperref rompa las referencias de longtable
\hypersetup{
    bookmarksnumbered = true,
    pdftitle = {Informe 2 - Modelamiento UML del sistema BeeTracer},
    pdfsubject = {ICI-3242 Modelamiento de Software - PUCV}
}
